\documentclass[11pt,a4paper]{article}
\usepackage{curso}
\begin{document}
\portada{Lección}{5.6}{No idealidades analógicas}

\noindent\textbf{Unidad 5.}\quad\textbf{Objetivos:} cuantificar el efecto de la cuantización y del desajuste sobre la distribución muestreada; no
linealidad y ruido de disparo; mitigación (calibración, redondeo estocástico, cuantizadores conservadores).
\textbf{Prerrequisitos:} 4.2, 4.5, 0.3 (KL).

\section{Medir el error en la distribución}
La métrica relevante no es el error de cada acoplamiento sino $\mathrm{KL}(p_{\text{ideal}}\|p_{\text{chip}})$. Ejemplo: Ising de 10 espines con
$J_{ij}\sim\N(0;0{,}5^2)$, $h_i\sim\N(0;0{,}3^2)$, enumerado exactamente (1024 estados) \evS.

\section{Cuantización y desajuste}
\begin{center}\begin{tabular}{lccccc}\toprule
bits de $J$ & 2 & 3 & 4 & 6 & 8\\\midrule
KL & 1,82 & 0,21 & 0,070 & 0,002 & $<10^{-3}$\\\bottomrule\end{tabular}\qquad
\begin{tabular}{lcccc}\toprule
$\sigma$ (desajuste de $h$) & 0,02 & 0,05 & 0,1 & 0,2\\\midrule
KL & 0,001 & 0,007 & 0,030 & 0,125\\\bottomrule\end{tabular}\end{center}
Para errores pequeños, $\mathrm{KL}\approx\tfrac{\beta^2}2\delta h^\top\mathrm{Cov}(s)\,\delta h\propto\sigma^2$: por eso se calibra cada dispositivo.

\section{No linealidad y ruido de disparo}
La suma del campo local no es perfectamente lineal y la curva de transferencia de un comparador con ruido de disparo es de tipo probit
($\Phi(\cdot)$), cercana pero distinta de una sigmoide. El ruido de disparo ($S_I=2qI$) no es térmico y cambia la temperatura efectiva $\beta_{\text{ef}}$.

\section{Mitigación}
\begin{definicion}[ (redondeo estocástico)]
$q(x)=\lfloor x/\Delta\rfloor\Delta+\Delta\,\mathbb 1[u<\{x/\Delta\}]$, $u\sim U(0,1)$: $\mathbb E[q(x)]=x$ (sin sesgo).
\end{definicion}
Promediando $K$ redondeos independientes a 3 bits, el error medio $|\hat J-J|$ baja de $0{,}10$ ($K=1$) a $0{,}04$ ($K=4$) y $0{,}023$ ($K=16$), mientras el
redondeo determinista queda en $0{,}088$. Cuantizadores conservadores: niveles elegidos para representar bien los acoplamientos grandes. En la salida hacen
falta $\ge8$ bits para no degradar los gradientes (Unidad 8). Diseñar un chip es un presupuesto de errores.

\begin{aplicacion}
Elegir bits desde una KL objetivo ($\ge6$ bits para KL $<10^{-2}$ aquí); calibrar; redondeo estocástico multiplexado en el tiempo; entrenamiento consciente del
hardware (lección 8.7).
\end{aplicacion}

\begin{error}
\begin{itemize}[nosep]
\item Juzgar por error de $J$ en vez de KL. \quad Suponer que el sesgo determinista se promedia.
\item Ignorar el ruido de disparo. \quad Calibrar una sola vez (deriva con la temperatura).
\end{itemize}
\end{error}

\begin{thebibliography}{9}\small
\bibitem[31]{aifer} M. Aifer et al., Error mitigation for thermodynamic computing, 2024.
\bibitem[48]{jel} A. Jelinčič et al., An efficient probabilistic hardware architecture for diffusion-like models, 2026.
\bibitem[41]{don} K. Donatella et al., Scalable thermodynamic second-order optimization, 2025.
\end{thebibliography}
\end{document}
